
%\documentclass[landscape]{slides}
\documentclass{beamer}
\usepackage{beamerfoils}
\usetheme{Pittsburgh}
%\usepackage[accumulated]{beamerseminar}
                                % remove ``accumulated'' option
                                % for original behaviour
%\usepackage{beamerthemeclassic}
\usepackage{ulem}
%\usepackage{beamerthemesplit}

%\usepackage{beamerthemeshadow}

%\usepackage{beamerthemeplain}
%\usepackage{times}
\usepackage{colortbl}

%\usepackage{lslide}
\def\emptyline{\vspace{10pt}}
\usepackage{amssymb}
%\usepackage{french}
\usepackage[latin1]{inputenc}
%\usepackage[T1]{fontenc}
\usepackage[]{color}
%\usepackage[usenames,dvipsnames]{color}
%\usepackage[pdftex]{color}
%\usepackage[]{colorlslide,epsfig}
%\usepackage{amsfonts,stmaryrd}
\usepackage{epsfig}

\usepackage{graphicx,array}
\usepackage{amsmath}%,amsthm}
\usepackage{hyperref}
\usepackage{pifont}

%\renewcommand{\frame}{\end{slide}   \begin{slide}}
%\newcommand{\frametitle}{\centerline}
 

%\usepackage{beamerthemeshadow}
\usepackage{graphicx}
%\usepackage{beamerthemesplit}



%%%%%%%%%%%%%%%%%% ENVIRONEMENT %%%%%%%%%%%%%%%%%%%%%%%%%%%%%%%%%%%%%%%%%%%%%%%%%
\newtheorem{deft}{D�finition}
%\newtheorem{theorem}{Th�or�me}
\newtheorem{proposition}{Proposition}
\newtheorem{exemple}{Exemple}
\newtheorem{postulat}{Postulat}


%%%%%%%%%%%%%%%%%% COULEURS %%%%%%%%%%%%%%%%%%%%%%%%%%%%%%%%%%%%%%%%%%%%%%%%%
%\definecolor{coltitle}{rgb}{0.1, 0.1, 0.543}
%\definecolor{colitemi}{rgb}{0, 0.75, 0.99}
%\definecolor{colexemp}{rgb}{0.6, 0.4, 0}
%\definecolor{colmult}{rgb}{0, 1, 0}
%\definecolor{gris}{rgb}{0.8, 0.8, 0.8}
\definecolor{encre}{rgb}{0.2, 0, 0.4}
\definecolor{lightblue}{rgb}{0.3,0.35,5}
%\definecolor{grisbleu}{rgb}{0.8, 0.8, 1}
%\definecolor{cadre}{rgb}{0.7, 0.7, 0.7}
%\definecolor{cadresubsec}{rgb}{0, 0.3, 0.3}
%\definecolor{cadreombre}{rgb}{0.8, 0.8, 0.8}
%\definecolor{colsection}{rgb}{0.4, 0.2, 0.2}

%\definecolor{bleur}{rgb}{0.3, 0, 0.6}
%\definecolor{bleuv}{rgb}{0, 0.3, 0.6}
%\definecolor{rougeb}{rgb}{0.6, 0, 0.3}
%\definecolor{rougev}{rgb}{0.6, 0.3, 0}
%\definecolor{rougealert}{rgb}{1., 0.2, 0}
%\definecolor{vertr}{rgb}{0.3, 0.6, 0}
%\definecolor{vertb}{rgb}{0, 0.6, 0.3}
%\definecolor{blanc}{rgb}{1, 1, 1}

\newcommand{\red}[1]{\only<#1>{\color{red}}}
\newcommand{\green}[1]{\only<#1>{\color{green}}}
\newcommand{\lra}{\leftrightarrow}
\newcommand{\Fd}{\mathbb{F}}
\newcommand{\enlight}[1]{\textit{\color{blue}#1}}
\newcommand{\blue}[1]{\textit{\color{blue}#1}}

%%   for the headline etc.
\renewcommand\normalcolor{\color{coltitle}}


%%   for the background
%\pagecolor{grisbleu}
%%   for the text
\color{encre}

%%%%%%%%%%%%%%% voir le style ....
 \beamertemplatenavigationsymbolsempty
\beamertemplatefootempty
%\beamertemplatefootframenumber
\setbeamertemplate{footline}[frame number]


\title{S�curit� informatique - pr�sentation et r�f�rences}

\AtBeginSection[]
{
\begin{frame}
\frametitle{Plan}
\tableofcontents[currentsection,sectionstyle=show/shaded,subsectionstyle=show/show/hide]
\end{frame}
}

\AtBeginSubsection[]
{
\begin{frame}
\frametitle{Plan}
\tableofcontents[currentsubsection,sectionstyle=show/shaded,subsectionstyle=show/shaded/hide]
\end{frame}
}
%%%%%%%%%%%%%%%%%%%%%%%%%%%%%%%%%%%%%%%%%%%%%%%%%%%%%%%%%%%%%%%%%%%%%%%%%%%%%%%%%%

\begin{document}

\maketitle



\begin{frame}

\frametitle{ \color{red} \Large Avertissement!}

\begin{itemize}\large  \color{red}
\item[\color{red}$\bullet$] Le but de ce cours est de vous �veillez aux �ventuels risques de s�curit� informatique.
\item[\color{red}$\bullet$] Le n'est  surtout pas de vous apprendre � devenir des pirates informatiques.
\item[\color{red}$\bullet$] Utilisez les techniques vues ici \textbf{ toujours } avec beaucoup de pr�cautions, uniquement dans un cadre d'apprentissage.
\end{itemize}

\end{frame}


\begin{frame}
  \frametitle{La s�curit� informatique }
  \begin{center}
  \includegraphics[width=8cm]{diagramme-securite.pdf}
\end{center}
\end{frame}


\begin{frame}
  \frametitle{Plan du cours}

  \begin{itemize}
  %\item {\color{blue} Chapitre 1.} Cryptographie,
  \item {\color{blue} Chapitre 1.} G�n�ralit� sur les malwares,
  \item {\color{blue} Chapitre 2.} Etude de virus par la pratique,
  \item {\color{blue} Chapitre 3.} Ver et m�thodes d'intrusion,
  \item {\color{blue} Chapitre 4.} S�curit� dans les r�seaux.
\end{itemize}

\end{frame}

\begin{frame}
\frametitle{R�f�rences}

\begin{center}
\begin{tabular}{cc}
 \textit{Les virus informatiques :}  & \textit{S�curit� informatique :} \\
\textit{th�orie, pratique et applications}  & \textit{Cours et exercices corrig�s.} \\
Eric Filiol. Springer. &  G. Avoine, P. Junod, P. Oechslin, Vuibert \\
\includegraphics[width=1.5cm]{couverture-filiol-virus-informatique.jpg} & \includegraphics[width=1.5cm]{securite-info-avoine.jpg} \\
& \\
& \\
\textit{Techniques de hacking.} &  \textit{Programmation syst�me sous Linux.} \\
 Jon Erickson. Pearson. &  Christophe Blaess. Eyrolles. \\
\includegraphics[width=1.5cm]{technique-de-hacking.jpg} & \includegraphics[width=1.5cm]{livre-progsys-41.jpeg} \\
\end{tabular}
\end{center}

\end{frame}

\end{document}
